\section{Проверка и сопровождение}
\label{sec:testing}

\subsection{Существующие тесты}
Из корня проекта:
\begin{lstlisting}[language=bash]
bash tests/run_tests.sh
python3 tests/cli_tests.py build-backend/PharmDelivery
\end{lstlisting}

\code{tests/run_tests.sh} создаёт временный каталог и удаляет его после
завершения. Компиляция выполняется с C++17, предупреждениями
\code{-Wall -Wextra -Wpedantic -Werror} и оптимизацией.
Компилятор можно задать через \code{CXX}.

\begin{tabularx}{\textwidth}{@{}L{0.30\textwidth}Y@{}}
\toprule
Набор & Основные проверки \\
\midrule
\code{tests/core_tests.cpp} & Полная, частичная и нулевая выдача; несколько
партий; скидки и порог крупной покупки; сроки годности и уценка; генерация;
спрос от наценки; приоритет уценённого; поставки; финансовый баланс. \\
\code{tests/simulation_tests.cpp} & Все восемь сочетаний границ $N=10/25$,
$M=3/9$, $K=15/35$; реальная нагрузка; очереди и опоздания; расписание;
повторные эксперименты; пустой склад; CSV; денежный и количественный баланс;
порядок вызовов. \\
\code{tests/cli_tests.py} & Справка и ошибочный ввод; параметры Qt-приложения;
читаемый JSON; соответствие сумм реальным доставкам; заданные скидки;
полная выдача при достаточном CSV-складе; повторяемость полного JSON. \\
\bottomrule
\end{tabularx}

Проверка ошибок памяти и неопределённого поведения:
\begin{lstlisting}[language=bash]
ASAN_OPTIONS=detect_leaks=0 SANITIZE=1 bash tests/run_tests.sh
\end{lstlisting}
Флаг \code{SANITIZE=1} добавляет AddressSanitizer и UndefinedBehaviorSanitizer.
\code{detect_leaks=0} отключает только проверку утечек, что используется
в средах, где LeakSanitizer недоступен из-за ptrace. Этот режим не заменяет
обычные проверки и не проверяет Qt-точку входа.

\subsection{Проверки при подготовке документации}
На 10 октября 2026 года текущие исходники были собраны через qmake
с Qt 5.15.19. В отдельной рабочей директории выполнены существующие
проверки, без изменения исходников и тестов проекта:
\begin{itemize}
  \item модульный набор: \code{Core checks passed: 1768};
  \item интеграционный набор: \code{Simulation checks passed: 76261};
  \item командная строка: \code{CLI checks passed: 1822}.
\end{itemize}
Число проверок характеризует конкретный запуск; оно может изменяться
при обновлении генерации или тестов. Режим санитайзеров при подготовке
этой документации не запускался.

\subsection{Диагностика типичных ситуаций}
\begin{description}
  \item[В первый день нет доставок.] Это ожидаемо: модель не создаёт
  вчерашних заказов до старта. Сегодняшняя подготовка доставляется завтра.
  \item[Не выполнен минимум 7 покупок.] Проверьте \code{delivered_orders}
  относительно $7M$. Недостаток спроса или товара не компенсируется
  искусственными заказами.
  \item[Возникла очередь или опоздания.] Сопоставьте новые запросы
  с вместимостью $15M$. Посмотрите \code{pending_orders} и
  \code{late_deliveries}; увеличение числа курьеров меняет вместимость.
  \item[Запрос был выполнен частично.] Проверьте \code{requested},
  \code{supplied} и срок годности партий на день предполагаемой доставки.
  Оставшаяся часть запроса автоматически не повторяется.
  \item[Цена партии уменьшилась, а закупочная стоимость прежняя.]
  Уценка меняет продажную базу, но учёт себестоимости сохраняет исходную цену.
  \item[Прибыль отличается от выручки минус закупки.]
  К денежному потоку добавляется закупочная стоимость оставшегося и
  подготовленного к доставке товара; см. раздел~\ref{sec:model}.
  \item[После CSV отсутствует часть лекарств.]
  Файл заменяет начальный склад целиком. Отсутствующие индексы имеют нулевой
  запас до поступления пополнения.
  \item[Ошибка последовательности модели.] Выполняйте четыре фазы дня
  в установленном порядке и читайте согласованный результат после завершения.
\end{description}

\subsection{Предположения и ограничения экспериментов}
Интенсивность обычного спроса, параметры выбора позиций, повышенный вес
уценённого товара, объём пополнения и интервалы доставки являются явными
правилами учебной модели. Они не выводятся из внешней статистики.
Выводы об эффективности сценариев относятся к этим предположениям.

Поток спроса непосредственно не зависит от числа курьеров, но вся модель
использует один генератор. Изменения наличия товаров, очереди, количества
покупателей и поставок могут сдвигать дальнейшую случайную последовательность.
Для сравнений фиксируйте параметры, \code{seed}, CSV и версию исходников;
при исследовании случайной вариативности используйте несколько значений seed.

Отдельные потоки, сохранение и восстановление состояния, отмена заказов,
повтор недовыданных позиций и продолжение за пределами $N$ дней не реализованы.
Открытые числовые поля требуют корректной инициализации, а прямое изменение
глобальных контейнеров требует соблюдения внутренних инвариантов.

\subsection{Обновление документации}
При изменении \code{SimulationParameters} обновите таблицу аргументов,
диапазоны и описание JSON-параметров. При изменении дневного движка сверяйте
порядок функций, признание выручки и состояние очередей. Изменение
\code{write_json_report} требует проверки всех таблиц формата и примеров.

Основные источники правил: \code{back/simulation.h},
\code{internal/simulation_engine.cpp}, \code{internal/simulation_cli.cpp},
\code{internal/supply_manager.cpp}, \code{internal/model_support.cpp},
\code{store.cpp}, \code{order.cpp}, \code{data_generator.cpp} и \code{tests/}.
Сборка самой документации описана в \code{docs/README.md}.
